\section{Details of Agent and Evaluation Protocols}

\subsection{Model and Execution Configuration}
\label{app:baseline_methods}
\textsc{OS-Omni} runs agents through execution specifications that bind together a task list, model identifier, observation mode, action mode, and history window. Model endpoints, provider-specific identifiers, authentication variables, and generation parameters are specified through a shared model registry, implemented in the release as \texttt{configs/models.yaml}. This separation keeps model selection independent from task initialization and evaluation, so the same task specifications can be reused across providers without changing the benchmark contract.

Unless otherwise stated, the main protocol uses screenshot-only observations and the platform action backend selected by the execution specification. The default history window contains the three most recent interaction steps. Primary tasks use the long-horizon limits declared in the task specification, with \texttt{limits.max\_steps:100} and \texttt{limits.step\_timeout\_seconds:120} as the released default. The step timeout bounds one model and action turn rather than the whole episode.

\begin{table}[ht]
\centering
\caption{Default execution parameters used in \textsc{OS-Omni} experiments.}
\label{tab:agent_execution_config}
\small
\begin{tabular}{p{0.32\linewidth}p{0.56\linewidth}}
\toprule
Parameter & Setting \\
\midrule
Observation mode & \texttt{screenshot\_only} \\
Action mode & Platform-specific controller selected by the execution specification \\
History window & 3 recent steps \\
Maximum steps & 100 per primary task \\
Step timeout & 120 seconds per step \\

\bottomrule
\end{tabular}
\end{table}



\subsection{Prompt Details}
\label{app:prompting_details}
Prompts are constructed from the task instruction, the platform-specific control interface, the current observation, and a bounded history of recent steps. In the main screenshot-only protocol, the agent is instructed to ground actions in the visible screen, use the coordinate frame of the current screenshot, avoid hidden application state or evaluator internals, and return one executable action at a time. Before the action, the agent provides a short reflection on the current screen and recent progress.

Across platforms, the agent follows the same high-level contract: it first gives a brief reflection, then returns an executable action block or one of the episode-control blocks \texttt{WAIT}, \texttt{DONE}, and \texttt{FAIL}. The platform runner validates the action block and dispatches it through the corresponding desktop, mobile, or browser controller. This design matches the shared action-space abstraction in the main paper while preserving the native execution channel used by each runtime.

Representative screenshot-only system prompts for the OS-specific runners are shown in Figures~\ref{fig:linux_prompt}--\ref{fig:ios_prompt}. Web tasks reuse the same benchmark-facing contract with browser-specific runtime instructions.
\clearpage
\begin{figure}[p]
\centering
\includegraphics[width=0.95\linewidth]{\detokenize{figures/Appendix_C(system_prompt)/Linux_prompt.pdf}}
\caption{System prompt used for Ubuntu/Linux screenshot-only runs.}
\label{fig:linux_prompt}
\end{figure}

\begin{figure}[p]
\centering
\includegraphics[width=0.95\linewidth]{\detokenize{figures/Appendix_C(system_prompt)/Windows_Prompt.pdf}}
\caption{System prompt used for Windows screenshot-only runs.}
\label{fig:windows_prompt}
\end{figure}

\begin{figure}[p]
\centering
\includegraphics[width=0.95\linewidth]{\detokenize{figures/Appendix_C(system_prompt)/mac_Prompt.pdf}}
\caption{System prompt used for macOS screenshot-only runs.}
\label{fig:mac_prompt}
\end{figure}

\begin{figure}[p]
\centering
\includegraphics[width=0.85\linewidth]{\detokenize{figures/Appendix_C(system_prompt)/Android_Prompt.pdf}}
\caption{System prompt used for Android screenshot-only runs.}
\label{fig:android_prompt}
\end{figure}

\begin{figure}[p]
\centering
\includegraphics[width=0.95\linewidth]{\detokenize{figures/Appendix_C(system_prompt)/iOS_Prompt.pdf}}
\caption{System prompt used for iOS screenshot-only runs.}
\label{fig:ios_prompt}
\end{figure}

\clearpage


\subsection{Action Parsing and Episode Execution}
\label{app:action_execution_protocol}

At each step, the runner sends the task instruction, the current screenshot observation, and the bounded recent history to the model. The model response is parsed into two parts: a short reflection and an executable action block. The reflection is stored in the trajectory for later analysis but is not executed. The action block is passed to the platform runner, which validates the format and dispatches the corresponding mouse, keyboard, touch, or device-control operation.

For desktop platforms, the action block is executed under a restricted PyAutoGUI-style sandbox or the corresponding desktop runtime controller. The sandbox permits direct mouse and keyboard control, short waits, and simple control flow, while disallowing hidden automation channels such as shell commands, screenshot APIs, image-location APIs, clipboard-based shortcuts, or application-internal scripting. For mobile platforms, the runner maps the same action intent to ADB-backed Android control or Appium/XCUITest-backed iOS control. For web tasks, browser-facing actions are dispatched through the runtime backend associated with the selected browser and host platform. This keeps the model-facing interaction protocol unified while allowing each environment to use its native execution backend.


After each executed action, the environment waits for the interface to settle, captures the next screenshot, and appends the response, action, controller status, and observation artifact to the trajectory. An episode terminates when the agent emits \texttt{DONE} or \texttt{FAIL}, when the task reaches the maximum step budget, or when a runtime condition such as a step timeout prevents further execution. The evaluator is invoked only after termination and checks the final environment state rather than the intermediate action sequence.

\subsection{Set-of-Mark Implementation Details}
\label{app:som_details}
Set-of-Mark (SoM) overlays are treated as optional visual grounding aids, not as a default benchmark observation. A SoM run annotates candidate regions on the screenshot and exposes mark identifiers to the agent. This can help isolate whether a model fails because it cannot localize a target or because it cannot plan the workflow after localization. Since SoM changes the visual input, SoM results should be reported as an ablation rather than mixed with screenshot-only results.

\subsection{Result Reporting Protocol}
\label{app:full_results}
Cross-platform result tables are valid only when platform lanes are run under a declared task inventory, model registry, observation protocol, action mode, and failure-accounting policy. The report should distinguish task failure from invalid runs caused by runtime setup, model-service errors, or evaluator exceptions. This convention follows the execution-based protocol in the main paper: a model succeeds only when the final environment state satisfies the task evaluator, while infrastructure failures are accounted for separately.

\paragraph{Recommended reporting fields.}
For each model and platform lane, reports include task count, success rate, invalid-run rate, average steps for successful tasks, timeout rate, observation mode, action mode, and the execution specification used for the run. Aggregate tables separate screenshot-only, accessibility-tree-only, combined-observation, and SoM settings. This prevents a visually grounded run from being compared against a metadata-assisted run as if the agent received the same information.

